\begin{lemma}[Principal-value cancellation of the Gaussian Riesz kernel on constants]
Fix $n \geq 1$, $j \in \{1,\ldots,n\}$ and $x \in \mathbb{R}^n$.  For $s > 0$ write
\[
  r = e^{-s} \in (0,1), \qquad \sigma = 1 - e^{-2s} = 1 - r^2 \in (0,1),
  \qquad
  p_s(x,y) = \frac{1}{\pi^{n/2}\sigma^{n/2}}\exp\Bigl(-\frac{|rx-y|^2}{\sigma}\Bigr) ,
\]
and let
\[
  D_j^x p_s(x,y) = \sqrt2\,r\,\frac{y_j - r x_j}{\sigma}\,p_s(x,y)
\]
be the energy-normalized $x$-derivative of the Lebesgue kernel, as evaluated in
\noderef{MehlerKernelXDerivative} (the displayed right-hand side is the
closed-form expression that node proves $2^{-1/2}\partial_{x_j}p_s(x,y)$ equals,
and it is that explicit expression, not a derivative, that occurs below).  Let
$K_{\gamma,j}$ be the normalized Gaussian Riesz kernel of
\noderef{GaussRieszKernel}.  Then:
\begin{enumerate}
\item for every $s > 0$ the function $y \mapsto D_j^xp_s(x,y)$ is Lebesgue
integrable on $\mathbb{R}^n$ and
\[
  \int_{\mathbb{R}^n} D_j^xp_s(x,y)\,dy = 0 ;
\]
\item for every $\varepsilon > 0$ the function $y \mapsto K_{\gamma,j}(x,y)$ is
Lebesgue integrable on the exterior region $\{y : |x-y| > \varepsilon\}$; and
\item
\[
  \lim_{\varepsilon\downarrow0}
    \int_{|x-y|>\varepsilon} K_{\gamma,j}(x,y)\,dy = 0 ,
\]
the limit being taken along $\varepsilon \to 0$ with $\varepsilon > 0$.
\end{enumerate}

Clause (3) is the reference paper's appendix display asserting that the
truncated kernel integrals of the constant function $1$ tend to zero, i.e.\ that
no diagonal multiplication term is left undetermined by the principal-value
realization; clauses (1) and (2) are the two well-definedness facts the
appendix establishes on the way, and both are used again downstream, clause (1)
in the subtracted-input step and clause (2) to know that the truncated integrals
of the appendix's final splitting are genuine Lebesgue integrals.  Every
constant occurring in the proof is allowed to depend on $n$, on $j$ and on $x$;
as the appendix header instructs, these constants serve only to identify two
realizations of one operator and no claim of uniformity in $n$ is made or used.
\end{lemma}
\begin{proof}
Throughout, $\lambda$ denotes Lebesgue measure on $\mathbb{R}^n$ and $n \geq 1$,
$j$ and $x$ are fixed.  Fix also, for the whole proof, the abbreviations: for
$s > 0$,
\[
  r = e^{-s} \in (0,1),\quad \sigma = 1-r^2 \in (0,1),\quad d = (1-r)x,
\]
and, for $z \in \mathbb{R}^n$,
\[
  q_\sigma(z) = (\pi\sigma)^{-n/2}e^{-|z|^2/\sigma},
  \qquad
  V_{j,\sigma}(z) = \frac{z_j}{\sigma}\,q_\sigma(z) .
\]
These are the appendix's letters; they are local notation for this proof and not
new objects.  Note $\sigma = 1-e^{-2s}$ because $r^2 = e^{-2s}$, and that
$0 < r < 1$ and $0 < \sigma < 1$ because $s > 0$.

\emph{Step 0: the kernel derivative in the appendix's letters.}
For $y \in \mathbb{R}^n$ put $z = y - x$.  Then
\[
  z + d = (y-x) + (1-r)x = y - rx ,
\]
so $|z+d| = |y-rx| = |rx-y|$ and therefore
$q_\sigma(z+d) = (\pi\sigma)^{-n/2}e^{-|rx-y|^2/\sigma} = p_s(x,y)$, while
$(z+d)_j = y_j - rx_j$.  Consequently
\begin{equation*}
  D_j^xp_s(x,y)
    = \sqrt2\,r\,\frac{(z+d)_j}{\sigma}\,q_\sigma(z+d)
    = \sqrt2\,r\,V_{j,\sigma}(z+d),
  \qquad z = y-x . \tag{0}
\end{equation*}
In particular $y \mapsto D_j^xp_s(x,y)$ is continuous, hence Borel measurable.

\emph{Step 1: Gaussian scaling identities.}
For $c > 0$ the dilation $w \mapsto cw$ of $\mathbb{R}^n$ multiplies Lebesgue
measure by $c^n$, $\lambda(cB) = c^n\lambda(B)$ for Borel $B$; hence for every
Borel $h : \mathbb{R}^n \to [0,+\infty]$,
\begin{equation*}
  \int_{\mathbb{R}^n} h(z)\,dz = c^n\int_{\mathbb{R}^n} h(cw)\,dw . \tag{1}
\end{equation*}
By \noderef{GaussMeasure} the measure $\gamma_n$ is Lebesgue measure with
density $\pi^{-n/2}e^{-|w|^2}$, and by \noderef{GaussIsProb} it is a probability
measure; hence
\begin{equation*}
  \int_{\mathbb{R}^n} \pi^{-n/2}e^{-|w|^2}\,dw = 1 . \tag{2}
\end{equation*}
Put
\[
  a_n = \int_{\mathbb{R}^n} |w_j|\,d\gamma_n(w),
  \qquad
  b_n = \int_{\mathbb{R}^n} |w_j|\,|w|\,d\gamma_n(w) .
\]
Both are finite: $|w_j| \leq |w| \leq 1 + |w|^2$ and $|w_j||w| \leq |w|^2$, and
$1+|w|^2$ and $|w|^2$ are polynomials in the coordinates of $w$, so finiteness
follows from the Gaussian exponential moments item of the shared context (with
$t = 0$) together with \noderef{GaussMeasure}.

Apply $(1)$ with $c = \sqrt\sigma$ and the three choices
$h = q_\sigma$, $h = |z_j|q_\sigma$ and $h = |z_j|\,|z|\,q_\sigma$, using
$q_\sigma(\sqrt\sigma w) = (\pi\sigma)^{-n/2}e^{-|w|^2}$,
$|(\sqrt\sigma w)_j| = \sqrt\sigma|w_j|$ and $|\sqrt\sigma w| = \sqrt\sigma|w|$,
and then $(2)$:
\begin{align}
  \int_{\mathbb{R}^n} q_\sigma(z)\,dz &= \sigma^{n/2}(\pi\sigma)^{-n/2}\!\!
     \int_{\mathbb{R}^n}\!\! e^{-|w|^2}dw = 1 , \tag{3a}\\
  \int_{\mathbb{R}^n} |z_j|\,q_\sigma(z)\,dz &= \sqrt\sigma\,a_n , \tag{3b}\\
  \int_{\mathbb{R}^n} |z_j|\,|z|\,q_\sigma(z)\,dz &= \sigma\,b_n . \tag{3c}
\end{align}
From $(3b)$,
\begin{equation*}
  \int_{\mathbb{R}^n}|V_{j,\sigma}(z)|\,dz
    = \frac1\sigma\int_{\mathbb{R}^n}|z_j|q_\sigma(z)\,dz
    = \frac{a_n}{\sqrt\sigma} < \infty , \tag{4}
\end{equation*}
so $V_{j,\sigma}$ is Lebesgue integrable.  By $(0)$ and the translation
invariance of Lebesgue measure (shared context), applied first to $y = x+z$ and
then to $z \mapsto z+d$,
\begin{equation*}
  \int_{\mathbb{R}^n}\bigl|D_j^xp_s(x,y)\bigr|\,dy
    = \sqrt2\,r\int_{\mathbb{R}^n}|V_{j,\sigma}(z+d)|\,dz
    = \sqrt2\,r\,\frac{a_n}{\sqrt\sigma} < \infty . \tag{5}
\end{equation*}
In particular $y \mapsto D_j^xp_s(x,y)$ is Lebesgue integrable on
$\mathbb{R}^n$, which is the integrability assertion of clause (1).

\emph{Step 2: the gradient and translation bounds for $V_{j,\sigma}$.}
The function $z \mapsto -|z|^2/\sigma$ is a polynomial in the coordinates with
gradient $-2z/\sigma$, so by the chain rule $q_\sigma$ is smooth with
$\nabla q_\sigma(z) = -\frac{2z}{\sigma}q_\sigma(z)$, and therefore
$V_{j,\sigma}$ is smooth with
\[
  \nabla V_{j,\sigma}(z)
    = \frac{q_\sigma(z)}{\sigma}\,e_j - \frac{2z_j}{\sigma^2}\,q_\sigma(z)\,z ,
\]
$e_j$ being the $j$th standard basis vector.  By the triangle inequality and
$|e_j| = 1$,
\[
  |\nabla V_{j,\sigma}(z)|
    \leq \frac{q_\sigma(z)}{\sigma} + \frac{2|z_j|\,|z|}{\sigma^2}q_\sigma(z) ,
\]
so integrating and using $(3a)$ and $(3c)$,
\begin{equation*}
  \int_{\mathbb{R}^n}|\nabla V_{j,\sigma}(z)|\,dz
    \leq \frac1\sigma + \frac{2}{\sigma^2}\,\sigma\,b_n
    = \frac{1+2b_n}{\sigma}
    =: \frac{C^{(1)}_n}{\sigma} . \tag{6}
\end{equation*}
Next, for any $e \in \mathbb{R}^n$ and any $z$, the function
$t \mapsto V_{j,\sigma}(z+te)$ is smooth on $[0,1]$, so by the fundamental
theorem of calculus
\[
  V_{j,\sigma}(z+e) - V_{j,\sigma}(z)
    = \int_0^1 \nabla V_{j,\sigma}(z+te)\cdot e\,dt ,
\]
whence $|V_{j,\sigma}(z+e)-V_{j,\sigma}(z)| \leq |e|\int_0^1
|\nabla V_{j,\sigma}(z+te)|\,dt$ by Cauchy--Schwarz (shared context) applied
pointwise in $t$.  The function $(z,t) \mapsto |\nabla V_{j,\sigma}(z+te)|$ is
continuous, hence product-measurable and nonnegative, so Tonelli (shared
context) and the translation invariance of $\lambda$ for each fixed $t$ give
\begin{equation*}
  \int_{\mathbb{R}^n}\bigl|V_{j,\sigma}(z+e)-V_{j,\sigma}(z)\bigr|\,dz
    \leq |e|\int_0^1\!\!\int_{\mathbb{R}^n}|\nabla V_{j,\sigma}(z+te)|\,dz\,dt
    \leq \frac{C^{(1)}_n\,|e|}{\sigma} . \tag{7}
\end{equation*}

\emph{Step 3: oddness, and clause (1).}
Let $R_j : \mathbb{R}^n \to \mathbb{R}^n$ be the reflection in the $j$th
coordinate, $R_jz = z - 2z_je_j$; it is linear, satisfies $R_j\circ R_j =
\mathrm{id}$ and preserves the Euclidean norm.  Lebesgue measure on
$\mathbb{R}^n$ is the $n$-fold product of one-dimensional Lebesgue measure, and
$R_j$ acts as the reflection $v \mapsto -v$ in the $j$th factor and as the
identity in the others; one-dimensional Lebesgue measure is invariant under that
reflection (shared context), and a product of measure-preserving maps of the
factors preserves the product measure, so $\lambda$ is invariant under $R_j$.
Since $|R_jz| = |z|$ we have
$q_\sigma\circ R_j = q_\sigma$, while $(R_jz)_j = -z_j$; hence
\[
  V_{j,\sigma}\circ R_j = -V_{j,\sigma} .
\]
Let $A \subseteq \mathbb{R}^n$ be Borel with $R_j(A) = A$; both $A = \mathbb{R}^n$
and $A = \{|z| > \varepsilon\}$ for $\varepsilon > 0$ have this property, the
latter because $|R_jz| = |z|$.  Since $V_{j,\sigma}$ is integrable by $(4)$,
invariance of $\lambda$ under $R_j$ gives
\[
  \int_A V_{j,\sigma}(z)\,dz
    = \int_A V_{j,\sigma}(R_jz)\,dz
    = -\int_A V_{j,\sigma}(z)\,dz ,
\]
so
\begin{equation*}
  \int_{\mathbb{R}^n} V_{j,\sigma}\,dz = 0
  \qquad\text{and}\qquad
  \int_{|z|>\varepsilon} V_{j,\sigma}\,dz = 0
  \quad(\varepsilon>0) . \tag{8}
\end{equation*}
Now by $(0)$, the translation invariance of $\lambda$ (first $y = x+z$, then
$z \mapsto z+d$) and the first identity of $(8)$,
\[
  \int_{\mathbb{R}^n} D_j^xp_s(x,y)\,dy
    = \sqrt2\,r\int_{\mathbb{R}^n}V_{j,\sigma}(z+d)\,dz
    = \sqrt2\,r\int_{\mathbb{R}^n}V_{j,\sigma}(w)\,dw
    = 0 ,
\]
all the integrals being absolutely convergent by $(4)$ and $(5)$.  Together with
$(5)$ this is clause (1).

\emph{Step 4: the $\varepsilon$-independent majorant.}
For $\varepsilon > 0$ and $s > 0$ set
$\Omega_\varepsilon = \{y \in \mathbb{R}^n : |x-y| > \varepsilon\}$ and
\[
  I_{\varepsilon,s} = \int_{\Omega_\varepsilon} D_j^xp_s(x,y)\,dy ,
\]
which is a well-defined real number by $(5)$.  Translating $y = x+z$, so that
$|x-y| = |z|$, and using $(0)$ and then the second identity of $(8)$,
\[
  I_{\varepsilon,s}
    = \sqrt2\,r\int_{|z|>\varepsilon}V_{j,\sigma}(z+d)\,dz
    = \sqrt2\,r\int_{|z|>\varepsilon}
        \bigl[V_{j,\sigma}(z+d)-V_{j,\sigma}(z)\bigr]\,dz ,
\]
both integrands being integrable by $(4)$ and $(7)$.  Hence by $(7)$ with
$e = d$,
\[
  |I_{\varepsilon,s}| \leq \sqrt2\,r\,\frac{C^{(1)}_n|d|}{\sigma} .
\]
Since $|d| = (1-r)|x|$ and $\sigma = 1-r^2 = (1-r)(1+r)$ with $1-r > 0$,
\[
  \frac{r|d|}{\sigma} = \frac{r(1-r)|x|}{(1-r)(1+r)}
    = |x|\,\frac{e^{-s}}{1+e^{-s}} \leq |x|\,e^{-s} ,
\]
which is the appendix's displayed majorant.  Therefore, with
$C^{(2)}_n = \sqrt2\,C^{(1)}_n$,
\begin{equation*}
  |I_{\varepsilon,s}| \leq C^{(2)}_n\,|x|\,e^{-s}
  \qquad\text{for every }\varepsilon>0\text{ and every }s>0 . \tag{9}
\end{equation*}
Moreover, for each fixed $s > 0$,
\begin{equation*}
  \lim_{\varepsilon\downarrow0} I_{\varepsilon,s} = 0 . \tag{10}
\end{equation*}
Indeed, by clause (1) already proved,
$|I_{\varepsilon,s}| = \bigl|\int_{\mathbb{R}^n}D_j^xp_s(x,y)\,dy
 - \int_{|x-y|\leq\varepsilon}D_j^xp_s(x,y)\,dy\bigr|
 = \bigl|\int_{|x-y|\leq\varepsilon}D_j^xp_s(x,y)\,dy\bigr|
 \leq \int_{\mathbb{R}^n}\mathbf 1_{\{|x-y|\leq\varepsilon\}}
   |D_j^xp_s(x,y)|\,dy$.
Given any sequence $\varepsilon_k \downarrow 0$ of positive reals, the
integrands $\mathbf 1_{\{|x-y|\leq\varepsilon_k\}}|D_j^xp_s(x,y)|$ converge
pointwise to $\mathbf 1_{\{y = x\}}|D_j^xp_s(x,y)|$, which vanishes
$\lambda$-almost everywhere since $\{x\}$ is Lebesgue-null, and they are
dominated by the integrable function $|D_j^xp_s(x,\cdot)|$ of $(5)$; dominated
convergence (shared context) gives $I_{\varepsilon_k,s} \to 0$.  As the sequence
was arbitrary and $\mathbb{R}$ is a metric space, so that convergence along
$\varepsilon \downarrow 0$ is equivalent to convergence along every such
sequence, $(10)$ follows.

\emph{Step 5: $K_{\gamma,j}$ as a time integral.}
We claim that for every $y \in \mathbb{R}^n$,
\begin{equation*}
  K_{\gamma,j}(x,y)
    = \frac{1}{\sqrt\pi}\int_0^\infty \frac{D_j^xp_s(x,y)}{\sqrt s}\,ds ,
  \tag{11}
\end{equation*}
with the convention --- the one \noderef{GaussRieszKernel} uses for its own
defining integral --- that an integral of a non-integrable integrand is $0$;
moreover the integrand on the right is Lebesgue integrable on $(0,\infty)$ if
and only if the integrand in the definition of $K_{\gamma,j}(x,y)$ is Lebesgue
integrable on $(0,1)$.

To see this, expand the right-hand integrand.  With $r = e^{-s}$ and
$\sigma = 1-r^2$,
\[
  \frac{1}{\sqrt\pi}\,\frac{D_j^xp_s(x,y)}{\sqrt s}
    = \frac{\sqrt2}{\pi^{(n+1)/2}}\;
      \frac{r\,(y_j-rx_j)}{\sqrt s\,(1-r^2)^{n/2+1}}\,
      \exp\Bigl(-\frac{|rx-y|^2}{1-r^2}\Bigr) ,
\]
where we used $p_s(x,y) = \pi^{-n/2}(1-r^2)^{-n/2}e^{-|rx-y|^2/(1-r^2)}$ and
$\pi^{-1/2}\pi^{-n/2} = \pi^{-(n+1)/2}$.  Let
\[
  h(\varrho) = \frac{y_j - \varrho x_j}{(1-\varrho^2)^{n/2+1}\sqrt{-\log \varrho}}
     \exp\Bigl(-\frac{|\varrho x-y|^2}{1-\varrho^2}\Bigr),
  \qquad \varrho \in (0,1) ,
\]
so that by \noderef{GaussRieszKernel},
$K_{\gamma,j}(x,y) = \sqrt2\,\pi^{-(n+1)/2}\int_{(0,1)}h(\varrho)\,d\varrho$.
The map $\psi(s) = e^{-s}$ is a continuously differentiable bijection of
$(0,\infty)$ onto $(0,1)$ with $\psi'(s) = -e^{-s} \neq 0$ and inverse
$\varrho \mapsto -\log\varrho$; in particular $-\log\psi(s) = s$, so
$\sqrt{-\log\psi(s)} = \sqrt s$, and $|\psi'(s)| = e^{-s}$.  By the
one-dimensional change-of-variables formula for such a bijection, which inserts
the factor $|\psi'|$, applied first to $|h|$ (a nonnegative Borel function, both
sides then being elements of $[0,+\infty]$) and then, when those are finite, to
$h$ itself,
\[
  \int_{(0,1)}h(\varrho)\,d\varrho = \int_{(0,\infty)}h(e^{-s})\,e^{-s}\,ds ,
\]
with the left integrand Lebesgue integrable exactly when the right one is.  A
direct substitution gives
\[
  h(e^{-s})\,e^{-s}
    = \frac{r\,(y_j-rx_j)}{(1-r^2)^{n/2+1}\sqrt s}
      \exp\Bigl(-\frac{|rx-y|^2}{1-r^2}\Bigr) ,
\]
which is exactly $\pi^{(n+1)/2}2^{-1/2}$ times the displayed right-hand
integrand above.  Multiplying through by $\sqrt2\,\pi^{-(n+1)/2}$ yields $(11)$
together with the stated equivalence of integrability; and when neither
integrand is integrable both sides of $(11)$ are $0$ by the convention.

\emph{Step 6: Fubini on $\Omega_\varepsilon \times (0,\infty)$, and clause (2).}
Fix $\varepsilon > 0$ and put
$F(y,s) = s^{-1/2}D_j^xp_s(x,y)$ for $(y,s) \in \Omega_\varepsilon\times(0,\infty)$.
It is continuous in $(y,s)$, hence product-measurable.  We bound
$\int_{\Omega_\varepsilon}|F(y,s)|\,dy$ in two regimes.  Put
\[
  s_0 = \min\Bigl(1,\ \frac{\varepsilon}{2(1+|x|)}\Bigr) > 0 .
\]

\emph{Small $s$.}  Let $0 < s \leq s_0$.  Since $1-e^{-s} \leq s$ we get
$|d| = (1-r)|x| \leq s|x| \leq \varepsilon|x|/(2(1+|x|)) \leq \varepsilon/2$.
Translating $y = x+z$ and then $w = z+d$, the region
$\{|z| > \varepsilon\}$ becomes $\{|w-d| > \varepsilon\}$, and
$|w-d| > \varepsilon$ with $|d| \leq \varepsilon/2$ forces
$|w| \geq |w-d| - |d| > \varepsilon/2$.  Hence, using $(0)$ and
$|w_j| \leq |w|$,
\[
  \int_{\Omega_\varepsilon}|D_j^xp_s(x,y)|\,dy
    \leq \sqrt2\,r\int_{|w|>\varepsilon/2}\frac{|w|}{\sigma}q_\sigma(w)\,dw .
\]
By $(1)$ with $c = \sqrt\sigma$ this equals
$\sqrt2\,\frac{r}{\sqrt\sigma}\int_{|u|>T}|u|\,\pi^{-n/2}e^{-|u|^2}\,du$ with
$T = \varepsilon/(2\sqrt\sigma)$.  For $|u| > T$ we have
$e^{-|u|^2} = e^{-|u|^2/2}e^{-|u|^2/2} \leq e^{-T^2/2}e^{-|u|^2/2}$, so
\[
  \int_{|u|>T}|u|\,\pi^{-n/2}e^{-|u|^2}\,du \leq \kappa_n\,e^{-T^2/2},
  \qquad
  \kappa_n := \pi^{-n/2}\int_{\mathbb{R}^n}|u|\,e^{-|u|^2/2}\,du ,
\]
and $\kappa_n < \infty$: by $(1)$ with $c = \sqrt2$,
$\int |u|e^{-|u|^2/2}du = 2^{n/2}\sqrt2\int|v|e^{-|v|^2}dv$, finite by the
Gaussian exponential moments item of the shared context (with $t=0$ and the
polynomial $1+|v|^2 \geq |v|$).  Finally $e^{-t} \leq 1/t$ for $t > 0$ with
$t = \varepsilon^2/(8\sigma)$ gives
$e^{-T^2/2} = e^{-\varepsilon^2/(8\sigma)} \leq 8\sigma/\varepsilon^2$, so,
using $r \leq 1$ and $\sigma \leq 1$,
\begin{equation*}
  \int_{\Omega_\varepsilon}|D_j^xp_s(x,y)|\,dy
    \leq \frac{8\sqrt2\,\kappa_n\,r\sqrt\sigma}{\varepsilon^2}
    \leq \frac{8\sqrt2\,\kappa_n}{\varepsilon^2}
  \qquad (0 < s \leq s_0) . \tag{12}
\end{equation*}

\emph{Large $s$.}  Let $s \geq s_0$.  Then
$\sigma = 1-e^{-2s} \geq 1-e^{-2s_0} > 0$, so with
$c_0 = (1-e^{-2s_0})^{-1/2} < \infty$, $(5)$ gives
\begin{equation*}
  \int_{\Omega_\varepsilon}|D_j^xp_s(x,y)|\,dy
    \leq \int_{\mathbb{R}^n}|D_j^xp_s(x,y)|\,dy
    \leq \sqrt2\,a_n\,c_0\,e^{-s}
  \qquad (s \geq s_0) . \tag{13}
\end{equation*}

Combining, and using $\int_0^{s_0}s^{-1/2}\,ds = 2\sqrt{s_0} < \infty$ for the
first regime and \noderef{SubordinationWeightIntegral} for the second,
\[
  \int_{(0,\infty)}\!\!\int_{\Omega_\varepsilon}|F(y,s)|\,dy\,ds
    \leq \frac{8\sqrt2\,\kappa_n}{\varepsilon^2}\cdot 2\sqrt{s_0}
      + \sqrt2\,a_n\,c_0\int_0^\infty \frac{e^{-s}}{\sqrt s}\,ds
    < \infty ,
\]
the iterated integral of the nonnegative product-measurable $|F|$ being the
product integral by Tonelli (shared context).  Hence $F$ is integrable on
$\Omega_\varepsilon\times(0,\infty)$, and Fubini (shared context) applies: for
$\lambda$-almost every $y \in \Omega_\varepsilon$ the slice
$s \mapsto F(y,s)$ is integrable on $(0,\infty)$, the function
$y \mapsto \int_0^\infty F(y,s)\,ds$ --- read with the same convention, so that
it is defined at \emph{every} $y$, with value $0$ at the exceptional ones --- is
Lebesgue integrable on $\Omega_\varepsilon$, and
\begin{equation*}
  \int_{\Omega_\varepsilon}\Bigl(\int_0^\infty F(y,s)\,ds\Bigr)dy
    = \int_0^\infty\Bigl(\int_{\Omega_\varepsilon}F(y,s)\,dy\Bigr)ds
    = \int_0^\infty \frac{I_{\varepsilon,s}}{\sqrt s}\,ds . \tag{14}
\end{equation*}
By $(11)$, $\int_0^\infty F(y,s)\,ds = \sqrt\pi\,K_{\gamma,j}(x,y)$ for every
$y$ --- for those $y$ at which the $s$-integrand is integrable this is $(11)$
itself, and at the remaining $y$ both sides are $0$ by the convention.  So
$y \mapsto \sqrt\pi K_{\gamma,j}(x,y)$ is Lebesgue integrable on
$\Omega_\varepsilon$, which is clause (2), and dividing $(14)$ by $\sqrt\pi$,
\begin{equation*}
  \int_{|x-y|>\varepsilon}K_{\gamma,j}(x,y)\,dy
    = \frac{1}{\sqrt\pi}\int_0^\infty \frac{I_{\varepsilon,s}}{\sqrt s}\,ds .
  \tag{15}
\end{equation*}

\emph{Step 7: dominated convergence in $\varepsilon$, and clause (3).}
By Fubini in Step 6 the function $s \mapsto I_{\varepsilon,s}$ is measurable on
$(0,\infty)$ for each $\varepsilon > 0$, and by $(9)$
\[
  \Bigl|\frac{I_{\varepsilon,s}}{\sqrt s}\Bigr|
    \leq C^{(2)}_n|x|\,\frac{e^{-s}}{\sqrt s}
  \qquad\text{for all }\varepsilon>0,\ s>0 ,
\]
the majorant being Lebesgue integrable on $(0,\infty)$ by
\noderef{SubordinationWeightIntegral}.  Let $\varepsilon_k \downarrow 0$ be any
sequence of positive reals.  By $(10)$, $I_{\varepsilon_k,s} \to 0$ for every
fixed $s > 0$, so dominated convergence (shared context, which is stated for
sequences) gives
$\int_0^\infty s^{-1/2}I_{\varepsilon_k,s}\,ds \to 0$.  Since the sequence was
arbitrary and $\mathbb{R}$ is a metric space, the limit along
$\varepsilon\downarrow0$ exists and equals $0$.  By $(15)$ this is clause (3).
\end{proof}
